%Version 3.1 December 2024
% See section 11 of the User Manual for version history
%
%%%%%%%%%%%%%%%%%%%%%%%%%%%%%%%%%%%%%%%%%%%%%%%%%%%%%%%%%%%%%%%%%%%%%%
%%                                                                 %%
%% Please do not use \input{...} to include other tex files.       %%
%% Submit your LaTeX manuscript as one .tex document.              %%
%%                                                                 %%
%% All additional figures and files should be attached             %%
%% separately and not embedded in the \TeX\ document itself.       %%
%%                                                                 %%
%%%%%%%%%%%%%%%%%%%%%%%%%%%%%%%%%%%%%%%%%%%%%%%%%%%%%%%%%%%%%%%%%%%%%

%%\documentclass[referee,sn-basic]{sn-jnl}% referee option is meant for double line spacing

%%=======================================================%%
%% to print line numbers in the margin use lineno option %%
%%=======================================================%%

%%\documentclass[lineno,pdflatex,sn-basic]{sn-jnl}% Basic Springer Nature Reference Style/Chemistry Reference Style

%%=========================================================================================%%
%% the documentclass is set to pdflatex as default. You can delete it if not appropriate.  %%
%%=========================================================================================%%

%%\documentclass[sn-basic]{sn-jnl}% Basic Springer Nature Reference Style/Chemistry Reference Style

%%Note: the following reference styles support Namedate and Numbered referencing. By default the style follows the most common style. To switch between the options you can add or remove “Numbered” in the optional parenthesis. 
%%The option is available for: sn-basic.bst, sn-chicago.bst%  
 
%%\documentclass[pdflatex,sn-nature]{sn-jnl}% Style for submissions to Nature Portfolio journals
%%\documentclass[pdflatex,sn-basic]{sn-jnl}% Basic Springer Nature Reference Style/Chemistry Reference Style
\documentclass[pdflatex,sn-mathphys-num]{sn-jnl}% Math and Physical Sciences Numbered Reference Style
%%\documentclass[pdflatex,sn-mathphys-ay]{sn-jnl}% Math and Physical Sciences Author Year Reference Style
%%\documentclass[pdflatex,sn-aps]{sn-jnl}% American Physical Society (APS) Reference Style
%%\documentclass[pdflatex,sn-vancouver-num]{sn-jnl}% Vancouver Numbered Reference Style
%%\documentclass[pdflatex,sn-vancouver-ay]{sn-jnl}% Vancouver Author Year Reference Style
%%\documentclass[pdflatex,sn-apa]{sn-jnl}% APA Reference Style
%%\documentclass[pdflatex,sn-chicago]{sn-jnl}% Chicago-based Humanities Reference Style

%%%% Standard Packages
%%<additional latex packages if required can be included here>

\usepackage{graphicx}%
\usepackage{multirow}%
\usepackage{amsmath,amssymb,amsfonts}%
\usepackage{amsthm}%
\usepackage{mathrsfs}%
\usepackage[title]{appendix}%
\usepackage{xcolor}%
\usepackage{textcomp}%
\usepackage{manyfoot}%
\usepackage{booktabs}%
\usepackage{algorithm}%
\usepackage{algorithmicx}%
\usepackage{algpseudocode}%
\usepackage{tcolorbox}
\usepackage{listings}%

\usepackage{changes}
    \definechangesauthor[name={Editor}, color=orange]{e}
	\definechangesauthor[name={Reviewer1}, color=blue]{r1}
	\definechangesauthor[name={Reviewer2}, color=red]{r3}
    \definechangesauthor[name={Reviewer2}, color=red]{r5}
    \definechangesauthor[name={Reviewer2}, color=red]{r10}
%%%%

%%%%%=============================================================================%%%%
%%%%  Remarks: This template is provided to aid authors with the preparation
%%%%  of original research articles intended for submission to journals published 
%%%%  by Springer Nature. The guidance has been prepared in partnership with 
%%%%  production teams to conform to Springer Nature technical requirements. 
%%%%  Editorial and presentation requirements differ among journal portfolios and 
%%%%  research disciplines. You may find sections in this template are irrelevant 
%%%%  to your work and are empowered to omit any such section if allowed by the 
%%%%  journal you intend to submit to. The submission guidelines and policies 
%%%%  of the journal take precedence. A detailed User Manual is available in the 
%%%%  template package for technical guidance.
%%%%%=============================================================================%%%%

%% as per the requirement new theorem styles can be included as shown below
\theoremstyle{thmstyleone}%
\newtheorem{theorem}{Theorem}%  meant for continuous numbers
%%\newtheorem{theorem}{Theorem}[section]% meant for sectionwise numbers
%% optional argument [theorem] produces theorem numbering sequence instead of independent numbers for Proposition
\newtheorem{proposition}[theorem]{Proposition}% 
%%\newtheorem{proposition}{Proposition}% to get separate numbers for theorem and proposition etc.

\theoremstyle{thmstyletwo}%
\newtheorem{example}{Example}%
\newtheorem{remark}{Remark}%

\theoremstyle{thmstylethree}%
\newtheorem{definition}{Definition}%

\raggedbottom
%%\unnumbered% uncomment this for unnumbered level heads

\begin{document}

\title[Adaptive Prompt Evolution for Continual Learning in Diabetic Retinopathy Detection]{Adaptive Prompt Evolution for Continual Learning in Diabetic Retinopathy Detection}

%%=============================================================%%
%% GivenName	-> \fnm{Joergen W.}
%% Particle	-> \spfx{van der} -> surname prefix
%% FamilyName	-> \sur{Ploeg}
%% Suffix	-> \sfx{IV}
%% \author*[1,2]{\fnm{Joergen W.} \spfx{van der} \sur{Ploeg} 
%%  \sfx{IV}}\email{iauthor@gmail.com}
%%=============================================================%%

\author*[1]{\fnm{Gusseppe} \sur{Bravo-Rocca}}\email{gusseppe.bravo@bsc.es}
\author[1]{\fnm{Peini} \sur{Liu}}\email{peini.liu@bsc.es}
\author[1,2]{\fnm{Jordi} \sur{Guitart}}\email{jordi.guitart@bsc.es}
\author[3]{\fnm{Ajay} \sur{Dholakia}}\email{adholakia@lenovo.com}
\author[3]{\fnm{David} \sur{Ellison}}\email{dellison@lenovo.com}
\author[4]{\fnm{Rodrigo M} \sur{Carrillo-Larco}}\email{rmcarri@emory.edu}

\affil[1]{\orgname{Barcelona Supercomputing Center}, \orgaddress{\city{Barcelona}, \country{Spain}}}
\affil[2]{\orgname{Universitat Politècnica de Catalunya}, \orgaddress{\city{Barcelona}, \country{Spain}}}
\affil[3]{\orgname{Lenovo Infrastructure Solutions Group}, \orgaddress{\city{Morrisville}, \state{NC}, \country{USA}}}
\affil[4]{\orgname{Emory University}, \orgaddress{\city{Atlanta}, \state{GA}, \country{USA}}}
%%==================================%%
%% Sample for unstructured abstract %%
%%==================================%%

\abstract{We introduce Adaptive Prompt Evolution (APE), a framework that autonomously optimizes text prompts for continual learning (CL) in medical imaging applications. Building upon task-agnostic domain-incremental learning, APE addresses static prompt engineering limitations through intelligent feedback loops that evolve templates and medical descriptions using iterative performance assessment. Our approach replaces manual prompt crafting with a local medical language model (Med-Gemma 4B) that generates and refines prompt candidates based on zero-shot clustering performance, maintaining complete offline operation for healthcare compliance. APE demonstrates remarkable efficiency, requiring only 50 evaluations to achieve F1-score of 0.9229, compared to exhaustive search requiring 110 evaluations for F1-score of 0.8921, representing 2.2× computational efficiency with 3.45\% performance improvement. Evaluated on diabetic retinopathy detection across three increasingly complex tasks, we integrate APE with established CL strategies including Elastic Weight Consolidation (EWC), Gradient Episodic Memory (GEM), and Learning Without Forgetting (LwF). On experiments across Multi-Layer Perceptron (MLP), Residual, and Attention architectures, APE outperforms baselines across all strategies. The framework maintains privacy-preserving experience replay using only embeddings while demonstrating robust generalization performance across tasks with minimal degradation.}



\keywords{Prompt Evolution, Continual Learning, Diabetic Retinopathy Detection}

%%\pacs[JEL Classification]{D8, H51}

%%\pacs[MSC Classification]{35A01, 65L10, 65L12, 65L20, 65L70}

\maketitle

\section{Introduction}
\label{sec:introduction}

Modern artificial intelligence systems require the ability to learn continuously from evolving data streams while preserving previously acquired knowledge \cite{Parisi2019}. This paradigm, known as Continual Learning (CL), presents particular challenges in medical imaging, where data privacy constraints and domain shifts complicate traditional learning approaches \cite{wang2024,Delange2021,Kirkpatrick2017,Robins1993}. Recent developments in vision-language models have enabled zero-shot learning capabilities through text prompts \cite{Radford2021}. However, current approaches in medical CL rely on static prompt engineering, where text descriptions are manually crafted or generated once and remain fixed throughout the learning process \cite{BravoRocca2025}. This static approach introduces several limitations: suboptimal prompts may compromise performance across diverse medical imaging conditions, manual prompt crafting requires domain expertise and is time-intensive, and fixed descriptions cannot adapt to evolving data distributions.

In medical imaging applications, these limitations are particularly relevant due to the variability of clinical data. For instance, diabetic retinopathy (DR) detection across different imaging devices or populations may benefit from distinct textual descriptions that capture domain-specific visual characteristics \cite{koh2022online}. Furthermore, reliance on external Large Language Model (LLM) APIs for prompt generation raises privacy and deployment concerns in healthcare environments where data sovereignty and offline operation are required \cite{kumar2018example,kumari2023deep,lenga2020continual}. To address these challenges, we introduce Adaptive Prompt Evolution (APE), a framework that autonomously discovers and optimizes text prompts for CL in medical imaging. APE implements an evolution loop that starts from basic descriptions and iteratively refines them through performance feedback. As shown in Figure \ref{fig:ape_overview}, the framework integrates Med-Gemma 4B \cite{medgemma-hf}, a specialized medical language model, to generate clinically relevant prompt candidates while maintaining complete local operation.

\begin{figure*}[htbp]
\centering
\includegraphics[width=1.0\linewidth]{ape.pdf}
\caption{APE framework for autonomous prompt optimization in medical CL. The system begins with multimodal inputs (image $x_i$, label $y_i$) and evolves from basic descriptions to optimal prompts through iterative feedback. Med-Gemma 4B generates candidate descriptions that are evaluated using zero-shot clustering in the CLIP embedding space. Best-performing prompts are selected based on clustering accuracy, enabling enhanced experience replay (ER) with privacy-preserving embedding storage. The optimized prompts generalize to subsequent tasks for multi-head classification without requiring labels.}
\label{fig:ape_overview}
\end{figure*}

Building upon our previous work on zero-shot clustering for CL \cite{BravoRocca2025}, APE transforms the approach from prompt usage to prompt learning. The system begins with elementary descriptions such as 'healthy' versus 'diseased' and autonomously evolves these into sophisticated medical terminology through iterative evaluation and selection. APE continues evolution beyond initial performance targets to optimize for cross-task generalization, addressing the challenge of knowledge generalization in medical CL scenarios \cite{khan2023looking,kuang2018stable,Zhang2024,Serra2018}. Our APE framework provides two technical contributions: autonomous prompt discovery through performance-driven evolution, and integration of local medical language model for privacy-preserving operation.

We evaluate APE through experiments on DR detection across three increasingly complex tasks \cite{aptos2019-blindness-detection}, integrating with established CL strategies including Elastic Weight Consolidation (EWC) \cite{Kirkpatrick2017}, Gradient Episodic Memory (GEM) \cite{LopezPaz2017}, and Learning Without Forgetting (LwF) \cite{Li2017}. Results demonstrate that APE achieves F1-score of 0.9229 versus 0.8921 for exhaustive search while requiring 2.2× fewer evaluations.

The main contributions of this paper are:
\begin{enumerate}
    \item A framework for autonomous prompt evolution that replaces static prompt engineering with adaptive learning for medical CL applications.
    \item Integration of Med-Gemma 4B for local prompt optimization that ensures privacy compliance while leveraging medical domain expertise.
    \item Integration and evaluation with multiple CL strategies and neural network architectures.
\end{enumerate}





%%%%%%%%%%%%%%%%%%%%%%%%%%%%%%%%%%%%%%%%%%%%%%%%%%%%%%%%%%%%%%%






\section{Related work}
\label{sec:related}

\noindent\textbf{LLMs for zero-shot learning}. Large Language Models (LLMs) have enhanced zero-shot learning capabilities through text-based instructions and prompts \cite{Brown2020}. In vision-language models like CLIP \cite{Radford2021}, LLMs generate textual descriptions that are converted into embeddings and compared with image embeddings for classification tasks. Our research extends this paradigm by using LLMs not just to create static descriptions, but to autonomously evolve and optimize these descriptions based on performance feedback. However, existing approaches rely on fixed prompts that limit adaptability to evolving medical imaging scenarios, motivating our autonomous evolution framework.

\noindent\textbf{Experience replay}. ER is rooted in emulating aspects of human memory processes, where past experiences are revisited to solidify learning \cite{Riemer18}. The canonical form involves interleaved training of new tasks with memory samples, seeking to approximate the joint distribution of tasks. Variants like Dark Experience Replay \cite{Buzzega20} have added complexity, employing distillation loss to enforce output consistency. Recent trends have incorporated dual-memory architectures, mirroring the interplay between fast and slow learning processes by maintaining two semantic memories \cite{Arani21}. While such architectures provide mechanisms to handle forgetting, the optimal way to structure and use these memories remains an open challenge.

\noindent\textbf{Privacy-preserving exemplars}. ER typically involves raw data samples from previous tasks for mitigating catastrophic forgetting. Our method enhances privacy by storing only embeddings of exemplars, not raw images. This modification ensures data privacy while maintaining DR effectiveness \cite{ICARL,DGR}. Through fine-tuning zero-shot clustering on training datasets, we refine exemplar selection, ensuring the memory buffer contains the most representative embeddings. This approach addresses privacy concerns in medical image analysis by storing embeddings instead of actual images, complying with privacy regulations and addressing security concerns \cite{ShokriShmatikov2015}.

\noindent\textbf{CLIP embeddings in CL}. Our methodology repurposes CLIP as an embedder within a CL framework, eschewing the common practice of fine-tuning CLIP on downstream tasks \cite{TiCCLIP}. This strategy retains the model's zero-shot learning capabilities while avoiding catastrophic forgetting inherent in direct fine-tuning scenarios. Through comparing LLM-generated descriptions, our zero-shot clustering process identifies optimal exemplars for memory storage, facilitating more effective learning across sequential tasks. APE extends this foundation through autonomous prompt optimization rather than static description selection.

\noindent\textbf{CL in medical imaging}. CL in medical imaging presents challenges due to privacy concerns and data distribution shifts. Recent work has explored continuous domain adaptation for healthcare applications \cite{venkataramani2018continuousdomainadaptationhealthcare}, addressing the evolving nature of medical data. Additionally, domain adaptation techniques have been applied to medical image segmentation tasks \cite{valindria2018domainadaptationmriorgan}, demonstrating the potential of transfer learning in this field. Our work builds upon these foundations, introducing autonomous prompt evolution that combines zero-shot learning with DR, specifically tailored to handle privacy and distribution shift issues in medical imaging scenarios.

\noindent\textbf{Prompt optimization and medical language models}. Recent developments in prompt optimization have introduced automated methods to enhance text-based instructions \cite{zhou2024systematicsurveypromptengineering}. However, these approaches often target general domains and single-task scenarios. In medical contexts, specialized language models like Med-PaLM \cite{singhal2022largelanguagemodelsclinical} demonstrate strong potential for domain-specific tasks \cite{singhal2023towardsexpertlevelmedical}, but they typically demand significant computational resources and external API access. Our APE framework overcomes these challenges by leveraging local deployment of Med-Gemma 4B \cite{medgemma-hf}, ensuring medical domain expertise, privacy compliance, and autonomous prompt evolution for CL.



%%%%%%%%%%%%%%%%%%%%%%%%%%%%%%%%%%%%%%%%%%%%%%%%%%%%%%%%%%%%%%%


\section{Background and problem formulation}
\label{sec:background}

\subsection{CL fundamentals}
In CL, a model learns from a sequence of tasks $\{\mathcal{T}_1, \mathcal{T}_2, \ldots, \mathcal{T}_n\}$ where each task represents a distinct data distribution. For task $\mathcal{T}_i$ with data $(\mathbf{x}_i, y_i)$, the objective is to minimize the cumulative loss across all tasks while preventing catastrophic forgetting:
\begin{equation}
\min_{\theta} \sum_{i=1}^{n} \mathcal{L}(\mathcal{H}(\mathbf{x}_i; \theta), y_i)
\end{equation}
where $\mathcal{H}$ represents the model with parameters $\theta$ and $\mathcal{L}$ denotes the loss function.

\subsection{TADILER framework and limitations}
Our work builds upon the TADILER framework \cite{BravoRocca2025}, which introduced zero-shot clustering for privacy-preserving CL in medical imaging. TADILER employed CLIP embeddings combined with LLM-generated prompts to achieve zero-shot classification, storing only embeddings rather than raw medical images for privacy compliance. 

However, TADILER suffered from significant limitations: it relied on exhaustive evaluation of 110 predefined prompt combinations, required external API dependencies for prompt generation, and most critically, used static prompts that could not adapt to domain-specific requirements or evolving data distributions. This static approach limits effectiveness across diverse medical imaging conditions and prevents optimization for specific clinical scenarios.

\subsection{Adaptive prompt optimization}
We extend TADILER to address autonomous prompt evolution in medical CL. Given medical images $\mathbf{x}_i$ and text prompts $\mathbf{p} = (t, d)$ consisting of templates $t$ and descriptions $d$, the zero-shot clustering function $\mathcal{C}(\mathbf{x}_i, \mathbf{p})$ assigns images to classes based on CLIP similarity. The objective is to find optimal prompts:
\begin{equation}
\mathbf{p}^* = \arg\max_{\mathbf{p}} \mathcal{F}(\mathcal{C}(\mathbf{x}_i, \mathbf{p}), y_i)
\end{equation}
where $\mathcal{F}$ represents the clustering evaluation metric (F1-score).

Unlike TADILER's static approach, adaptive prompt optimization autonomously discovers effective prompts through iterative refinement, enabling domain-specific adaptation while maintaining privacy-preserving characteristics. This approach presents several challenges: the search space for optimal prompts is combinatorially large, requiring intelligent exploration rather than exhaustive enumeration; medical LLMs must generate clinically relevant terminology while operating under local deployment constraints; and evolved prompts must generalize across medical imaging tasks without access to subsequent task labels. Our APE framework addresses these challenges through local medical LLM integration and performance-driven selection mechanisms.


%%%%%%%%%%%%%%%%%%%%%%%%%%%%%%%%%%%%%%%%%%%%%%%%%%%%%%%%%%%%%%%%


\section{Adaptive prompt evolution framework}
\label{sec:ape_framework}

\subsection{Framework overview}
APE operates through an iterative evolution loop that begins with elementary prompt descriptions and autonomously refines them based on zero-shot clustering performance. Unlike TADILER's exhaustive evaluation of 110 predefined combinations, APE intelligently explores the prompt space through guided evolution, requiring significantly fewer evaluations while achieving superior performance. The framework consists of four core components: (1) Med-Gemma 4B integration for medical prompt generation, (2) performance-driven candidate evaluation using CLIP embeddings, (3) evolution control mechanisms including stopping criteria, and (4) privacy-preserving DR with evolved prompts.

\subsection{Medical LLM integration and prompt generation}
We employ Med-Gemma 4B \footnote{https://huggingface.co/google/medgemma-4b-it}, a specialized medical LLM with 4 billion parameters, for autonomous prompt generation. This model has been pre-trained on a variety of medical data, including medical text relevant to radiology images, chest-x rays, histopathology patches, ophthalmology images and dermatology images. Med-Gemma 4B provides medical domain expertise while maintaining computational efficiency suitable for local deployment, addressing privacy concerns inherent in external API dependencies. The model operates entirely offline through local inference, eliminating data transmission to external services and configured with temperature 0.7 (more diverse and creative outputs) and maximum 300 tokens to balance creativity with coherence in medical terminology generation.

Given current best-performing prompts consisting of template $t$ and descriptions $d = \{d_{\text{healthy}}, d_{\text{diseased}}\}$, Med-Gemma 4B generates $k$ candidate combinations through structured prompting that leverages its medical training:

\begin{tcolorbox}[colback=blue!5!white, colframe=blue!75!black, title=\textbf{Med-Gemma 4B prompt}]
\begin{small}
\begin{verbatim}
You are a medical expert creating optimal text prompts for diabetic 
retinopathy classification using AI vision models.

CURRENT BEST (F1={current_f1:.3f}):
Template: "{current_template}"
Healthy: "{current_descriptions['No retinopathy']}"  
Diseased: "{current_descriptions['Retinopathy']}"

TASK: Generate {num_candidates} NEW template+description combinations 
that could beat F1={current_f1:.3f}

REQUIREMENTS:
- Include BOTH template and descriptions
- Keep ALL text SHORT (2-4 words max per description)  
- Use medical terminology that AI vision models understand
- Focus on visual features visible in fundus images
...
\end{verbatim}
\end{small}
\end{tcolorbox}

This structured approach ensures Med-Gemma 4B generates medically appropriate candidates while maintaining the brevity required for effective CLIP embedding generation. We found that longer descriptions do not improve scores because CLIP does not pay attention to details.

\subsection{Evolution algorithm}
APE implements an intelligent evolution loop that iteratively refines prompt candidates through performance-driven selection. The algorithm operates in iterative cycles: Med-Gemma 4B generates candidate prompt combinations based on current best-performing prompts, each candidate is evaluated using CLIP-based zero-shot clustering, and the best-performing candidate becomes the new baseline if it improves F1-score performance. The process uses early stopping when no improvement (step) occurs for consecutive iterations (to avoid local minimum).  Algorithm \ref{alg:ape_evolution} outlines the complete evolution process. Later, in Section \ref{sec:results_analysis}, we show how APE improves twice the best current score (2 steps) running 10 iterations in total.

\begin{algorithm}
\caption{Adaptive Prompt Evolution Algorithm}\label{alg:ape_evolution}
\begin{algorithmic}[1]
\Require Initial prompts $\mathbf{p}_0 = (t_0, d_0)$, target score $\tau$, max iterations $I_{\text{max}}$, patience $P$, candidates per iteration $k$, stop on target flag $\text{stop\_on\_target}$
\Ensure Evolved prompts $\mathbf{p}^*$, evolution tracking data $\mathcal{T}$
\State $\mathbf{p}_{\text{best}} \Leftarrow \mathbf{p}_0$
\State $f_{\text{best}} \Leftarrow \mathcal{F}(\mathcal{C}(\mathbf{x}, \mathbf{p}_0), y)$
\State $\text{patience\_count} \Leftarrow 0$
\State $\mathcal{T} \Leftarrow \{\}$ \Comment{Initialize tracking data}
\For{$i = 1$ to $I_{\text{max}}$}
    \State $\mathcal{P}_{\text{candidates}} \Leftarrow \text{GenerateCandidates}(\mathbf{p}_{\text{best}}, f_{\text{best}}, k)$
    \State $\text{improved} \Leftarrow \text{False}$
    \For{each candidate $\mathbf{p}_c \in \mathcal{P}_{\text{candidates}}$}
        \State $f_c \Leftarrow \mathcal{F}(\mathcal{C}(\mathbf{x}, \mathbf{p}_c), y)$
        \If{$f_c > f_{\text{best}}$}
            \State $\mathbf{p}_{\text{best}} \Leftarrow \mathbf{p}_c$
            \State $f_{\text{best}} \Leftarrow f_c$
            \State $\text{patience\_count} \Leftarrow 0$
            \State $\text{improved} \Leftarrow \text{True}$
            \State Record evolution step in $\mathcal{T}$
        \EndIf
    \EndFor
    \If{$\text{improved} = \text{False}$}
        \State $\text{patience\_count} \Leftarrow \text{patience\_count} + 1$
    \EndIf
    \If{$f_{\text{best}} \geq \tau$ \textbf{and} $\text{stop\_on\_target} = \text{True}$}
        \State \textbf{break} \Comment{Target achieved}
    \EndIf
    \If{$\text{patience\_count} \geq P$}
        \State \textbf{break} \Comment{Convergence detected}
    \EndIf
\EndFor
\State \Return $\mathbf{p}_{\text{best}}, \mathcal{T}$
\end{algorithmic}
\end{algorithm}

% Figure \ref{fig:learning_progression} illustrates the autonomous learning trajectory, showing how APE achieves substantial improvement (+0.308 F1-score) in just 2 evolution steps (2 successful improvements), requiring only 50 total evaluations compared to TADILER's 110.



%%%%%%%%%%%%%%%%%%%%%%%%%%%%%%%%%%%%%%%%%%%%%%%%%%%%%%%%%%%%%

\subsection{Performance evaluation and evolution control}
Each generated candidate undergoes evaluation through the same zero-shot clustering mechanism employed by TADILER, ensuring consistency with the established framework. Given candidate prompts $\mathbf{p}_c = (t_c, d_c)$, we compute text embeddings using CLIP and perform clustering through cosine similarity:

\begin{equation}
S_{ij} = \cos(\mathbf{X}_i, \mathbf{T}_j) = \frac{\mathbf{X}_i \cdot \mathbf{T}_j}{|\mathbf{X}_i| |\mathbf{T}_j|}
\end{equation}

where $\mathbf{X}_i$ represents image embeddings and $\mathbf{T}_j$ represents text embeddings derived from candidate prompts. Performance assessment uses F1-score due to class imbalance in medical datasets:

\begin{equation}
\mathcal{F} = \frac{2 \cdot \text{Precision} \cdot \text{Recall}}{\text{Precision} + \text{Recall}}
\end{equation}

APE implements multiple stopping mechanisms to balance computational efficiency with performance optimization. The framework employs patience-based early stopping when no improvement occurs for $P$ consecutive iterations, preventing unnecessary computation while allowing sufficient exploration. Additionally, APE supports optional target-based stopping when a specified performance threshold is achieved, though our implementation continues evolution beyond targets to optimize for generalization.

Figure \ref{fig:candidate_quality_distribution} demonstrates the evolution process across iterations, showing how APE intelligently navigates the prompt space to identify high-performing combinations. The violin plots reveal that while candidate quality varies significantly within iterations, the evolution process successfully identifies and selects superior prompts (marked with stars) that drive performance improvements.

\begin{figure}[htbp]
\centering
\includegraphics[width=1.0\linewidth]{candidate_quality_distribution.pdf}
\caption{Candidate quality distribution across iterations showing APE's exploration strategy. Violin plots display F1-score distributions of all generated candidates per iteration, with background shading indicating performance zones (excellent $\!> 0.9$, good $0.8\!-\!0.9$, poor $< 0.8$). Gold stars mark candidates selected for evolution, connected by the red trend line showing consistent improvement despite varying candidate quality within each iteration.}
\label{fig:candidate_quality_distribution}
\end{figure}

\subsection{Exploration efficiency and prompt analysis}
The exploration analysis reveals that APE effectively balances diversity and quality, generating medically relevant terminology while maintaining the brevity required for effective CLIP embedding generation. The framework discovers 24 unique templates and 47 distinct description pairs across 50 total search points, with only 4 selections driving the evolution process, contrasting with exhaustive methods that evaluate all possible combinations regardless of effectiveness.

Figure \ref{fig:search_space_exploration} illustrates APE's navigation through the prompt landscape, while Figure \ref{fig:prompt_component_analysis} provides detailed analysis of how different templates and descriptions contribute to overall effectiveness in diabetic retinopathy detection tasks. In Figure \ref{fig:prompt_component_analysis}, we observe that some descriptions generally perform well (yellow), while others achieve the best results when used in combination (green).



\begin{figure}[htbp]
\centering
\includegraphics[width=1.0\linewidth]{search_space_exploration.pdf}
\caption{APE exploration timeline showing evolutionary progression through prompt optimization. The scatter plot displays all candidate prompts across iterations, colored by F1-score performance. Red circles trace the evolution path from initial baseline (F1=0.615) to major improvement with "Fundus image reveals" template (F1=0.921), followed by description refinement from complex terms ("Normal retina structure", "Diabetic retinopathy features") to simplified versions ("Normal retina", "Diabetic retinopathy") achieving final performance (F1=0.923). Annotation boxes detail the specific changes driving performance improvements.}
\label{fig:search_space_exploration}
\end{figure}

\begin{figure}[htbp]
\centering
\includegraphics[width=1.0\linewidth]{prompt_component_analysis.pdf}
\caption{Prompt component performance analysis across template, healthy term, and disease term categories. The horizontal bar charts show average F1-scores for different prompt components, with occurrence counts (n=X) indicating frequency in the evolution process. Gold highlighting identifies components with the highest average performance across all combinations, while green highlighting with star patterns marks components from the best final evolved combination (F1=0.923).}
\label{fig:prompt_component_analysis}
\end{figure}

%%%%%%%%%%%%%%%%%%%%%%%%%%%%%%%%%%%%%%%%%%%%%%%%%%%%%%%%%%%%

\subsection{Integration with CL strategies}
APE seamlessly integrates with established CL algorithms while maintaining TADILER's privacy-preserving characteristics. The evolved prompts enhance DR effectiveness across multiple CL paradigms, including Elastic Weight Consolidation (EWC), Gradient Episodic Memory (GEM), Learning Without Forgetting (LwF), and Naive approaches. The integration process follows Algorithm \ref{alg:ape_continual_learning}, which extends the TADILER DR mechanism with autonomously evolved prompts. The key advantage lies in the superior quality of evolved prompts for zero-shot clustering, which directly improves exemplar selection for DR. Unlike TADILER's static descriptions, APE's evolved prompts capture nuanced medical terminology that better distinguishes between healthy and pathological retinal images.

\begin{algorithm}
\caption{APE-Enhanced CL Integration}\label{alg:ape_continual_learning}
\begin{algorithmic}[1]
\Require Task sequence $\{\mathcal{T}_1, \mathcal{T}_2, \ldots, \mathcal{T}_n\}$, CL strategy $\mathcal{S}$, buffer size $B$
\Ensure Trained model $\mathcal{H}$ with evolved prompts $\mathbf{p}^*$
\State $\mathbf{p}^* \Leftarrow \text{APE-Evolution}(\mathcal{T}_1)$ \Comment{Evolve prompts on first task}
\For{each task $\mathcal{T}_i$ in sequence}
    \If{$i = 1$}
        \State Train model $\mathcal{H}$ on $\mathcal{T}_1$ using strategy $\mathcal{S}$
    \Else
        \State $\mathcal{E}_{\text{replay}} \Leftarrow \text{GenerateExemplars}(\mathbf{p}^*, \text{PreviousTasks}, B)$
        \State Train model $\mathcal{H}$ on $\mathcal{T}_i \cup \mathcal{E}_{\text{replay}}$ using strategy $\mathcal{S}$
    \EndIf
    \State Update experience buffer with evolved prompts $\mathbf{p}^*$
\EndFor
\State \Return $\mathcal{H}, \mathbf{p}^*$
\end{algorithmic}
\end{algorithm}

As demonstrated in the experimental results (Section \ref{sec:results_analysis}), APE shows consistent performance improvements across all CL strategies, with particularly strong gains in challenging scenarios such as knowledge distillation (LwF) and memory-based approaches (GEM).

\subsection{Privacy-preserving deployment and efficiency}
APE maintains TADILER's privacy-preserving paradigm while enhancing exemplar quality through evolved prompts. The framework stores only CLIP embeddings derived from optimal prompt configurations, ensuring medical image privacy compliance while improving clustering accuracy for DR. Given evolved prompts $\mathbf{p}^* = (t^*, d^*)$, the system generates optimized text embeddings:

\begin{equation}
\mathbf{T}^*_j = \frac{\text{CLIP}_{\text{text}}(t^* + d^*_j)}{|\text{CLIP}_{\text{text}}(t^* + d^*_j)|}
\end{equation}

where $d^*_j$ represents the evolved descriptions for each class. The stratified sampling procedure from TADILER is enhanced through these evolved prompts, with exemplar selection based on similarity to evolved text embeddings:

\begin{equation}
\mathcal{E}_c = \text{TopK}\left(\left\{(\mathbf{x}_i, S(\mathbf{X}_i, \mathbf{T}^*_c)) : y_i = c\right\}, k\right)
\end{equation}

where $S(\mathbf{X}_i, \mathbf{T}^*_c)$ represents cosine similarity between image and evolved text embeddings for class $c$.

APE's local deployment architecture addresses critical healthcare requirements while providing computational efficiency advantages. The framework operates entirely offline using Med-Gemma 4B, eliminating external API dependencies and associated privacy concerns. The evolution process typically converges within 2-3 iterations, requiring approximately 9.2 minutes total optimization time, making APE practical for deployment scenarios where prompt optimization must occur locally while maintaining rapid adaptation capabilities.




%%%%%%%%%%%%%%%%%%%%%%%%%%%%%%%%%%%%%%%%%%%%%%%%%%%%%%%%%%%%%%%%


\section{Experimental methodology}
\label{sec:experimental_methodology}

We conduct comprehensive experiments to evaluate APE's effectiveness in medical CL scenarios, comparing autonomous prompt evolution against static approaches across multiple architectures and CL strategies. Our experimental design extends the TADILER's evaluation framework \cite{BravoRocca2025} with APE-specific protocols to assess prompt optimization efficiency, generalization capabilities, and integration effectiveness.

\subsection{Experimental setup}
Experiments were conducted on a CPU-based platform with dual Intel\textsuperscript{\textregistered} Xeon\textsuperscript{\textregistered} Platinum 8360Y CPUs operating at 2.40GHz with 256 GB RAM, running Ubuntu 22.04 LTS. The software stack includes Intel\textsuperscript{\textregistered} AI Analytics Toolkit, Avalanche 0.3.1 for CL, Intel\textsuperscript{\textregistered} Extension for PyTorch 1.12.100+cpu, and Med-Gemma 4B deployed through LM Studio for local inference. This configuration ensures reproducibility while maintaining offline operation essential for healthcare compliance.

We use the APTOS 2019 Blindness Detection dataset \cite{aptos2019-blindness-detection}, comprising 3,662 high-resolution retinal images partitioned across three sequential tasks: Task 0 (832 training, 208 testing), Task 1 (2,097 training, 525 testing), and Task 2 (2,929 training, 733 testing). Following TADILER's task design, we construct three sequential tasks with increasing complexity: Task 0 with uniform image quality representing ideal clinical conditions, Task 1 introducing lighting variations simulating different imaging environments, and Task 2 incorporating Gaussian noise addition emulating degraded imaging conditions. Figure \ref{fig:fundus_images} shows representative samples from each task.

\begin{figure}[htbp]
\centering
\includegraphics[width=1.0\linewidth]{fundus_images.pdf}
\caption{Fundus images from different tasks. Left: Task 0 - uniform image quality; Center: Task 1 - variation in lighting; Right: Task 2 - artificially added Gaussian noise to simulate challenging imaging conditions \cite{BravoRocca2025}.}
\label{fig:fundus_images}
\end{figure}

\subsection{Model architectures and CL strategies}
We employ three distinct architectures, as shown in Figure \ref{fig:architectures}, processing 768-dimensional CLIP embeddings with 256 hidden layer size and 2-class output for binary DR classification:

\begin{figure}[htbp]
\centering
\includegraphics[width=1.0\linewidth]{architectures.pdf}
\caption{Neural network architectures used in APE evaluation: (a) Multi-Layer Perceptron with single hidden layer, (b) Residual Network with skip connections across three ResidualBlocks, and (c) Attention Network with 4-head self-attention mechanism. All architectures process 768-dimensional CLIP embeddings for binary DR  classification.}
\label{fig:architectures}
\end{figure}

\begin{itemize}
    \item \textbf{Multi-Layer Perceptron (MLP)}: Simple architecture with single hidden layer (Linear → ReLU → Linear) serving as baseline for fundamental enhancement assessment.
    \item \textbf{Residual Network}: Three-layer residual architecture with skip connections (ResidualBlock × 3 → Linear) mitigating vanishing gradients crucial for multi-task performance.
    \item \textbf{Attention Network}: Self-attention mechanism with 4 attention heads (Linear → ReLU → SelfAttention → Linear) capturing complex dependencies relevant to medical feature detection.
\end{itemize}

All models employ Adam optimizer with learning rate 0.001 and CrossEntropyLoss criterion, configured as multi-task architectures using Avalanche's as\_multitask wrapper. We implement four established CL approaches using consistent hyperparameters: training batch size 200, evaluation batch size 200, and 2 epochs per task. Strategy-specific configurations include Naive (standard fine-tuning), EWC with $\lambda = 0.2$, LwF with $\alpha = 0.5$ and temperature $T = 0.2$, and GEM with 10 patterns per experience. Each strategy integrates with both random exemplar selection and APE-enhanced TADILER selection for direct comparison.

\subsection{APE optimization and evaluation protocol}
APE optimization employed 7 iterations with 5 candidates per iteration, patience threshold of 3 iterations for convergence detection, and target F1-score of 0.8921 matching TADILER's exhaustive search baseline. Med-Gemma 4B operates with temperature 0.7 and maximum 300 tokens for local inference. The evolution process begins with elementary prompts (template: "", healthy: "healthy", diseased: "diseased") and autonomously refines them through performance feedback, continuing beyond target achievement to optimize cross-task generalization.

We implement comprehensive validation through multiple random seeds [3, 51, 35] ensuring result reliability, and neighbor count exploration [15, 20, 25, 30, 50] assessing stratified sampling sensitivity. The experimental sequence follows [0, 1, 2] task progression with systematic evaluation after each task completion. Our evaluation employs F1-score for zero-shot clustering effectiveness and Average Mean Class Accuracy (AMCA) as primary CL metric:

\begin{equation}
AMCA = \frac{1}{T} \sum_{t=1}^{T} \left( \frac{1}{C} \sum_{c=1}^{C} a_{c,t} \right)
\end{equation}

where $T=3$ represents tasks and $C=2$ represents classes. We establish two primary baselines for comparison: 'Original' approach using random exemplar selection without prompt optimization, and TADILER's exhaustive search across 110 predefined prompt combinations. 


\subsection{Computational complexity analysis}
\label{sec:complexity}
\added[id=r1]{
APE's time complexity is determined by three nested operations. Let $I$ denote the
maximum number of evolution iterations, $k$ the number of candidates generated per
iteration, $N$ the number of data points in the task, and $D$ the CLIP embedding
dimension. Each candidate evaluation computes cosine similarity between $N$ image
embeddings and two text embeddings, giving $\mathcal{O}(N \cdot D)$ per evaluation.
LLM-guided candidate generation adds a constant per-candidate inference cost
$C_{\text{LLM}}$ that is independent of $N$. The total complexity of one APE evolution
run is therefore:
\begin{equation}
\mathcal{O}(I \cdot k \cdot (N \cdot D + C_{\text{LLM}}))
\end{equation}
In our experiments, $I_{\text{max}} = 10$, $k = 5$, $D = 768$, and $N$ ranges from
832 to 2{,}929 depending on the task. In contrast, TADILER's exhaustive search
evaluates all combinations from a fixed vocabulary:
$\mathcal{O}(|\mathcal{T}| \cdot |\mathcal{H}| \cdot |\mathcal{D}| \cdot N \cdot D)$,
where $|\mathcal{T}|$, $|\mathcal{H}|$, $|\mathcal{D}|$ are the template, healthy
term, and disease term vocabulary sizes (110 combinations total). APE's patience
mechanism terminates once no improvement is observed for $P$ consecutive iterations,
giving worst-case bound $\mathcal{O}(I_{\text{max}} \cdot k \cdot N \cdot D)$ and
empirical convergence at 50 total evaluations, a 2.2$\times$ reduction over the
exhaustive 110-evaluation baseline. Empirically, APE converges in approximately 9.2
minutes on a CPU-only platform (dual Intel Xeon Platinum 8360Y), consistent with
practical offline deployment in resource-constrained healthcare environments.
}
%%%%%%%%%%%%%%%%%%%%%%%%%%%%%%%%%%%%%%%%%%%%%%%%%%%%%%%%%%%


\section{Results and analysis}
\label{sec:results_analysis}

We present comprehensive experimental results demonstrating APE's effectiveness across multiple dimensions: autonomous prompt optimization performance, integration with CL strategies, computational efficiency gains, and generalization capabilities. Our analysis reveals that APE consistently outperforms both original baselines and static TADILER approaches while requiring significantly fewer computational resources. While APE demonstrates consistent superiority across most architecture-strategy combinations, some cases show comparable performance to TADILER, particularly in simpler architectures where model capacity may limit the benefits of prompt evolution (as in MLP vs. LwF from Table \ref{tab:amca_comparison_lwf_ewc}).

\subsection{Autonomous prompt evolution performance}
APE demonstrates remarkable autonomous learning capability, evolving from basic prompt descriptions to sophisticated medical terminology through performance-driven optimization. The evolution process achieves F1-score of 0.9229 starting from elementary prompts (F1=0.6145), representing a substantial +0.308 improvement in just 2 evolution steps (improvements over current best score). This autonomous discovery process required only 50 candidate evaluations compared to TADILER's exhaustive 110-evaluation approach, demonstrating 2.2× computational efficiency while achieving superior performance (+3.45\% improvement over exhaustive search baseline of F1=0.8921). Figures \ref{fig:learning_progression} and \ref{fig:computational_efficiency} illustrate these performance and efficiency gains.


\begin{figure}[htbp]
\centering
\includegraphics[width=1.0\linewidth]{learning_progression.pdf}
\caption{APE learning progression from basic to optimal prompts. The system evolves from elementary descriptions achieving F1=0.615 to sophisticated medical terminology achieving F1=0.923, representing +0.308 improvement in 552.9 seconds with only 50 candidate evaluations (10 iterations x 5 candidates, the best score was found at iteration 7). Two improvements (steps) were found in all the process.}
\label{fig:learning_progression}
\end{figure}


\begin{figure}[htbp]
\centering
\includegraphics[width=0.7\linewidth]{computational_efficiency.pdf}
\caption{Computational efficiency comparison between APE and exhaustive search approaches. APE achieves superior performance (F1=0.9229 vs. 0.8921) while requiring 2.2× fewer evaluations (50 vs. 110), demonstrating efficiency gains through intelligent prompt space exploration.}
\label{fig:computational_efficiency}
\end{figure}


The evolved prompts transition from basic descriptions (Template: "", Healthy: "healthy", Disease: "diseased") to clinically relevant terminology (Template: "Fundus image reveals", Healthy: "Normal retina", Disease: "Diabetic retinopathy"), showcasing Med-Gemma 4B's ability to generate medically appropriate language while maintaining the brevity required for effective CLIP embedding generation. Generalization experiments demonstrate consistent performance maintenance with F1-scores of 0.917 and 0.918 for Tasks 1 and 2 respectively, representing minimal degradation from the learning task performance of 0.923, as show in Figure \ref{fig:performance_analysis}.


% Figure \ref{fig:performance_analysis} demonstrates APE's superior performance and effective cross-task generalization. APE achieves F1-score of 0.923 during learning (Task 0) compared to exhaustive search baseline of 0.892, representing a +0.031 improvement. The transfer performance to Tasks 1 and 2 shows minimal degradation, maintaining F1-scores of 0.917 and 0.918 respectively (+0.027 improvements over exhaustive search), demonstrating effective prompt generalization without requiring re-optimization for each new task.

\begin{figure}[htbp]
\centering
\includegraphics[width=1.0\linewidth]{performance_analysis.pdf}
\caption{Performance comparison between APE and exhaustive search when learning and generalizing. APE demonstrates superior performance during learning (Task 0) and maintains effectiveness during generalization to Tasks 1 and 2, showing minimal performance degradation while consistently outperforming exhaustive search baselines.}
\label{fig:performance_analysis}
\end{figure}



\subsection{CL integration and architectural analysis}
APE demonstrates consistent superiority across all CL strategies and neural network architectures. Figure \ref{fig:performance_strategies} illustrates APE's performance advantage across EWC (0.958 vs. 0.923 TADILER), GEM (0.963 vs. 0.946 TADILER), LwF (0.959 vs. 0.901 TADILER), and Naive (0.955 vs. 0.921 TADILER) strategies, with results averaged across all three neural network architectures (MLP, Residual, and Attention). The performance improvements are particularly pronounced in challenging scenarios: LwF shows the largest improvement (+0.058 AMCA points), indicating APE's effectiveness in knowledge distillation scenarios, while GEM demonstrates robust performance with minimal variance.

\begin{figure}[htbp]
\centering
\includegraphics[width=1.0\linewidth]{performance_strategies.pdf}
\caption{Performance comparison across CL strategies showing Mean AMCA scores for Original (red), TADILER (orange), and APE (green) approaches. APE consistently outperforms both baselines across all strategies: EWC (+0.035 vs. TADILER), GEM (+0.017 vs. TADILER), LwF (+0.058 vs. TADILER), and Naive (+0.034 vs. TADILER). Error bars represent standard deviation across multiple runs and architectures.}
\label{fig:performance_strategies}
\end{figure}

Tables \ref{tab:amca_comparison_naive_gem} and \ref{tab:amca_comparison_lwf_ewc} provide comprehensive performance analysis across architectures and neighbor configurations. Key findings include:

%%%%%%%%%%%%%%%%% Table 1

\begin{table*}
\centering
\tiny
  {\caption{Comparison of Mean AMCA scores and forgetting (parentheses) for Naive and GEM strategies across different base models. Boldface
indicates superior performance among \textit{Original}, \textit{TADILER}, and \textit{APE} methods based on Mean AMCA. Averages are
computed across different NN values for each base model. TADILER values taken from \cite{BravoRocca2025}.}
  \label{tab:amca_comparison_naive_gem}
}
  {\begin{tabular}{llcccccc}
  
  \toprule
  \bfseries Model & \bfseries NN & \multicolumn{3}{c}{\bfseries Naive} & \multicolumn{3}{c}{\bfseries GEM} \\
  \cmidrule(lr){3-5} \cmidrule(lr){6-8}
 &  \bfseries  & \bfseries Original & \bfseries TADILER & \bfseries APE & \bfseries Original & \bfseries TADILER & \bfseries APE \\
  \midrule
\multirow{7}{*}{\textbf{Attention}} & 15 & 0.929(3.4) & 0.936(3.7) & \textbf{0.967(3.5)} & 0.965(3.4) & 0.967(3.5) & \textbf{0.968(1.3)} \\
 & 20 & 0.926(4.3) & 0.938(4.2) & \textbf{0.963(4.2)} & 0.965(3.2) & \textbf{0.968(4.2)} & 0.967(2.4) \\
 & 25 & 0.922(4.9) & 0.939(4.7) & \textbf{0.966(4.0)} & 0.965(3.2) & 0.967(3.3) & \textbf{0.968(3.1)} \\
 & 30 & 0.930(2.5) & 0.937(4.5) & \textbf{0.965(5.1)} & 0.961(3.4) & 0.967(4.2) & \textbf{0.967(2.8)} \\
 & 50 & 0.932(2.6) & 0.938(4.3) & \textbf{0.958(5.7)} & 0.960(3.8) & 0.965(4.1) & \textbf{0.966(3.1)} \\
  \cmidrule(lr){2-8}
 & \textbf{Avg} & 0.928(3.5) & 0.938(4.3) & \textbf{0.964(4.5)} & 0.963(3.4) & 0.967(3.9) & \textbf{0.967(2.6)} \\
  \midrule
\multirow{7}{*}{\textbf{Residual}} & 15 & 0.880(9.7) & 0.913(5.1) & \textbf{0.965(3.6)} & 0.934(2.8) & 0.941(2.6) & \textbf{0.970(4.5)} \\
 & 20 & 0.879(9.7) & 0.912(5.0) & \textbf{0.966(5.5)} & 0.935(2.9) & 0.942(2.8) & \textbf{0.968(1.8)} \\
 & 25 & 0.879(9.7) & 0.913(5.3) & \textbf{0.969(5.5)} & 0.938(2.8) & 0.942(1.7) & \textbf{0.968(2.2)} \\
 & 30 & 0.879(9.6) & 0.912(5.7) & \textbf{0.967(5.8)} & 0.937(2.7) & 0.940(3.9) & \textbf{0.967(4.0)} \\
 & 50 & 0.903(3.8) & 0.912(5.0) & \textbf{0.967(5.0)} & 0.935(2.6) & 0.941(3.1) & \textbf{0.969(2.9)} \\
  \cmidrule(lr){2-8}
 & \textbf{Avg} & 0.884(8.5) & 0.912(5.2) & \textbf{0.967(5.1)} & 0.936(2.8) & 0.941(2.8) & \textbf{0.968(3.1)} \\
  \midrule
\multirow{7}{*}{\textbf{MLP}} & 15 & 0.899(3.2) & 0.915(5.1) & \textbf{0.935(5.0)} & 0.923(3.8) & 0.931(5.1) & \textbf{0.952(5.4)} \\
 & 20 & 0.897(4.3) & 0.914(5.3) & \textbf{0.934(5.6)} & 0.923(3.7) & 0.931(5.1) & \textbf{0.952(5.5)} \\
 & 25 & 0.900(3.8) & 0.914(5.2) & \textbf{0.935(5.5)} & 0.922(4.0) & 0.931(5.1) & \textbf{0.953(5.6)} \\
 & 30 & 0.896(4.2) & 0.914(5.6) & \textbf{0.935(5.5)} & 0.923(3.8) & 0.930(5.2) & \textbf{0.951(5.5)} \\
 & 50 & 0.893(5.1) & 0.911(5.5) & \textbf{0.937(5.6)} & 0.921(3.9) & 0.927(5.3) & \textbf{0.952(5.5)} \\
  \cmidrule(lr){2-8}
 & \textbf{Avg} & 0.897(4.1) & 0.914(5.4) & \textbf{0.935(5.4)} & 0.922(3.9) & 0.930(5.1) & \textbf{0.952(5.5)} \\
  \midrule
  \bottomrule
  \end{tabular}}
\end{table*}


%%%%%%%%%%%% Table 2 

\begin{table*}
\centering
\tiny
  {\caption{Comparison of Mean AMCA scores and forgetting (parentheses) for LwF and EWC strategies across different base models. Boldface
indicates superior performance among \textit{Original}, \textit{TADILER}, and \textit{APE} methods based on Mean AMCA. Averages are
computed across different NN values for each base model. TADILER values taken from \cite{BravoRocca2025}.}
  \label{tab:amca_comparison_lwf_ewc}
}
  {\begin{tabular}{llcccccc}
  
  \toprule
  \bfseries Model & \bfseries NN & \multicolumn{3}{c}{\bfseries LwF} & \multicolumn{3}{c}{\bfseries EWC} \\
  \cmidrule(lr){3-5} \cmidrule(lr){6-8}
 &  \bfseries  & \bfseries Original & \bfseries TADILER & \bfseries APE & \bfseries Original & \bfseries TADILER & \bfseries APE \\
  \midrule
\multirow{7}{*}{\textbf{Attention}} & 15 & 0.793(5.3) & 0.823(4.8) & \textbf{0.969(5.1)} & 0.928(3.1) & 0.930(4.1) & \textbf{0.965(5.1)} \\
 & 20 & 0.793(5.3) & 0.816(5.8) & \textbf{0.966(4.7)} & 0.919(4.4) & 0.934(3.7) & \textbf{0.968(4.0)} \\
 & 25 & 0.789(5.4) & 0.826(5.5) & \textbf{0.969(5.2)} & 0.915(4.2) & 0.933(3.6) & \textbf{0.970(4.3)} \\
 & 30 & 0.792(5.5) & 0.825(5.3) & \textbf{0.967(5.2)} & 0.918(4.1) & 0.929(5.9) & \textbf{0.969(4.8)} \\
 & 50 & 0.791(5.0) & 0.824(5.7) & \textbf{0.967(5.5)} & 0.911(5.7) & 0.932(4.8) & \textbf{0.968(3.2)} \\
  \cmidrule(lr){2-8}
 & \textbf{Avg} & 0.792(5.3) & 0.823(5.4) & \textbf{0.968(5.1)} & 0.918(4.3) & 0.932(4.4) & \textbf{0.968(4.3)} \\
  \midrule
\multirow{7}{*}{\textbf{Residual}} & 15 & 0.928(4.5) & 0.940(4.6) & \textbf{0.969(4.9)} & 0.933(3.7) & 0.939(3.8) & \textbf{0.969(3.8)} \\
 & 20 & 0.930(4.5) & 0.934(4.5) & \textbf{0.969(4.6)} & 0.932(4.1) & 0.938(4.3) & \textbf{0.970(2.5)} \\
 & 25 & 0.928(4.4) & 0.933(4.6) & \textbf{0.968(5.1)} & 0.926(4.5) & 0.938(3.9) & \textbf{0.970(4.0)} \\
 & 30 & 0.928(4.5) & 0.938(4.4) & \textbf{0.969(4.6)} & 0.914(6.7) & 0.938(4.9) & \textbf{0.968(2.9)} \\
 & 50 & 0.929(4.5) & 0.936(4.8) & \textbf{0.968(5.0)} & 0.901(10.0) & 0.932(6.7) & \textbf{0.968(2.9)} \\
  \cmidrule(lr){2-8}
 & \textbf{Avg} & 0.929(4.5) & 0.936(4.6) & \textbf{0.969(4.8)} & 0.921(5.8) & 0.937(4.7) & \textbf{0.969(3.2)} \\
  \midrule
\multirow{7}{*}{\textbf{MLP}} & 15 & 0.940(4.8) & \textbf{0.942(4.8)} & 0.941(5.1) & 0.917(5.8) & 0.930(6.9) & \textbf{0.935(6.0)} \\
 & 20 & 0.940(4.5) & \textbf{0.941(5.0)} & 0.941(5.1) & 0.874(5.4) & 0.892(5.2) & \textbf{0.936(5.8)} \\
 & 25 & 0.937(5.0) & 0.944(5.0) & \textbf{0.944(5.0)} & 0.875(4.1) & 0.891(5.6) & \textbf{0.934(5.6)} \\
 & 30 & 0.937(5.3) & \textbf{0.944(4.8)} & 0.943(5.0) & 0.879(3.2) & 0.892(5.5) & \textbf{0.936(5.8)} \\
 & 50 & 0.936(5.3) & \textbf{0.948(4.9)} & 0.941(5.0) & 0.872(4.7) & 0.890(5.9) & \textbf{0.939(5.1)} \\
  \cmidrule(lr){2-8}
 & \textbf{Avg} & 0.938(5.0) & \textbf{0.944(4.9)} & 0.942(5.0) & 0.883(4.6) & 0.899(5.8) & \textbf{0.936(5.7)} \\
  \midrule
  \bottomrule
  \end{tabular}}
\end{table*}



\textbf{Attention architecture}: APE achieves exceptional performance across all strategies, with particularly strong results in Naive (0.964 average AMCA) and LwF (0.968 average AMCA) approaches. The attention mechanism benefits significantly from APE's enhanced prompt quality, with improvements ranging from +0.026 (GEM) to +0.145 (LwF) over static TADILER approaches.

\textbf{Residual architecture}: Shows the most dramatic improvements, with APE achieving 0.967-0.969 AMCA across all strategies compared to original baselines ranging from 0.884-0.929. This suggests that APE's enhanced exemplar selection particularly benefits architectures with skip connections, enabling better knowledge preservation through improved DR.

\textbf{MLP architecture}: Demonstrates substantial improvements in EWC (0.936 vs. 0.883 original) and competitive performance in other strategies. Notably, in LwF, the simpler MLP architecture shows comparable performance between TADILER and APE (0.944 vs. 0.942), suggesting that complex architectures benefit more from autonomous prompt evolution.

Figure \ref{fig:amca_forgetting} presents the trade-off between performance retention and catastrophic forgetting across all methods. APE consistently achieves the ideal combination of high AMCA scores with low forgetting rates, occupying the upper-left region representing optimal CL performance. This dual benefit stems from improved exemplar selection quality through evolved prompts, enabling more effective DR that preserves previous task knowledge while facilitating new task learning.

\begin{figure}[htbp]
\centering
\includegraphics[width=1.0\linewidth]{amca_forgetting.pdf}
\caption{AMCA versus forgetting rate trade-off analysis. APE (green triangles) consistently occupies the ideal region of high AMCA and low forgetting rates. The clear separation between APE and other methods (Original - red circles, TADILER - orange squares) demonstrates superior CL capability.}
\label{fig:amca_forgetting}
\end{figure}

\subsection{Computational efficiency and robustness analysis}
APE's autonomous evolution demonstrates superior computational efficiency compared to exhaustive search approaches. The optimization process converges rapidly, typically within 2-3 iterations, requiring approximately 9.2 minutes total optimization time. Figure \ref{fig:performance_evolution} illustrates APE's stability across different neighbor configurations, showing consistent performance maintenance regardless of stratified sampling parameters. This robustness reduces hyperparameter sensitivity, simplifying deployment configuration.

\begin{figure}[htbp]
\centering
\includegraphics[width=1.0\linewidth]{performance_evolution.pdf}
\caption{Performance evolution across neighbor count (KNN) values showing stability of Original (red), TADILER (orange), and APE (green) approaches. APE maintains consistently superior performance (0.957 AMCA) across all neighbor configurations (15-50), demonstrating robustness to hyperparameter variations. Shaded regions represent confidence intervals across multiple runs and architectures.}
\label{fig:performance_evolution}
\end{figure}

Results across three random seeds (3, 51, 35) demonstrate APE's consistent performance advantages with statistical significance (p $<$ 0.05) across all strategy comparisons. Effect sizes range from moderate (GEM: +0.017) to large (LwF: +0.058). The neighbor count exploration reveals APE's robustness to stratified sampling configuration, with performance variations remaining within ±0.003 AMCA points across all neighbor counts, demonstrating that APE's benefits are not dependent on precise hyperparameter tuning.

APE's local deployment architecture provides scalability advantages for medical AI deployment, eliminating external API dependencies while ensuring privacy compliance. Med-Gemma 4B's computational requirements (CPU is enough) enable deployment on standard hardware configurations, making APE accessible for resource-constrained healthcare environments. The autonomous evolution process demonstrates Med-Gemma 4B's capability to generate clinically relevant terminology without explicit medical supervision, with evolved descriptions capturing essential visual characteristics that enhance CLIP's embedding quality compared to generic terms. This emergence of domain-specific language validates the local medical LLM integration approach while maintaining effectiveness across varying imaging conditions.



%%%%%%%%%%%%%%%%%%%%%%%%%%%%%%%%%%%%%%%%%%%%%%%%%

\section{Limitations}
\label{sec:limitations}

While APE demonstrates consistent improvements across medical CL scenarios, several limitations warrant consideration for future development and broader deployment.

The current evaluation focuses exclusively on DR detection within a single medical imaging modality. Generalization to other medical domains such as radiology, pathology, or dermatology remains to be validated. Different medical specialties may require distinct prompt evolution strategies or domain-specific LLMs beyond Med-Gemma 4B's current training scope. APE's reliance on Med-Gemma 4B introduces computational overhead compared to static prompt approaches. While it enables standard hardware deployment, resource-constrained environments may find the local LLM deployment challenging. Additionally, the evolution process requires initial optimization time that may not be suitable for scenarios demanding immediate deployment. 

The autonomous evolution framework currently optimizes for binary classification tasks (healthy vs diseased). Extension to multi-class medical classification scenarios would require modifications to the prompt generation strategy and evaluation metrics. Medical diagnosis often involves multiple severity levels or different pathological conditions that may benefit from more complex prompt structures. 


CLIP's embedding quality inherently limits APE's performance ceiling. While evolved prompts improve clustering effectiveness, fundamental limitations in vision-language model understanding of medical imagery may constrain potential improvements. Medical images often contain subtle visual features that general-purpose vision models may not capture optimally. In addition, reliance on embedding-based representations presents explainability challenges that might require integration with attention visualization techniques or other tools that provide insights into the embedding space.

The evaluation dataset, while clinically relevant, represents a controlled experimental setting. Real-world medical imaging deployment involves additional complexities including varying image acquisition protocols, diverse patient populations, and regulatory compliance requirements that extend beyond the current experimental scope.

Finally, APE's stopping criteria and generalization strategy require task-specific tuning. While the framework demonstrates robustness across neighbor configurations, optimal patience thresholds and iteration limits may vary across different medical applications, potentially requiring domain expert involvement for optimal configuration.

\section{Conclusion}
\label{sec:conclusion}

We introduced Adaptive Prompt Evolution (APE), a framework that transforms static prompt engineering into autonomous optimization for medical CL. APE addresses limitations of existing approaches through local medical LLM integration, performance-driven evolution, and privacy-preserving operation essential for healthcare deployment. Our evaluation on DR detection across three increasingly complex tasks demonstrates APE's effectiveness across multiple dimensions. The autonomous evolution process achieves F1-score of 0.9229 versus exhaustive search baseline of 0.8921 while requiring 2.2× fewer evaluations. Integration with established CL strategies shows consistent improvements: EWC (+0.035 AMCA), GEM (+0.017 AMCA), LwF (+0.058 AMCA), and Naive (+0.034 AMCA) approaches. APE's evolution strategy produces prompts that generalize effectively across tasks without re-optimization, maintaining F1-scores above 0.917 when generalizing. This capability addresses practical deployment requirements where prompt adaptation for each new medical domain may not be feasible. The framework's local operation through Med-Gemma 4B eliminates external API dependencies, ensuring healthcare compliance while leveraging medical domain expertise. The experimental results across three neural network architectures and multiple random seeds demonstrate statistical significance and robustness of APE's improvements. Particularly notable are the substantial gains in challenging scenarios such as LwF integration and residual architecture deployment, indicating APE's effectiveness across diverse CL paradigms.
APE establishes a shift from manual prompt engineering to autonomous optimization in medical AI applications. The framework enables healthcare organizations to deploy CL systems that adapt to evolving data distributions while maintaining privacy compliance and operational efficiency. Future work should explore extension to multi-class medical classification, validation across additional medical domains, and integration with federated learning scenarios for multi-institutional deployment. The convergence of autonomous prompt evolution, local medical LLMs, and privacy-preserving CL represents a step toward practical medical AI systems that can adapt continuously while meeting healthcare deployment requirements.
\section*{Acknowledgements}

We thank Lenovo for providing the technical infrastructure to run the experiments in this paper. This work was partially supported by Lenovo and Intel\textsuperscript{\textregistered} as part of the Lenovo AI Innovators University Research program, by the Spanish Ministry of Science (MICINN), the Research State Agency (AEI) and European Regional Development Funds (ERDF/FEDER) under grant agreement PID2021-126248OB-I00, MCIN/AEI/10.13039/501100011033/FEDER, UE, and by the Generalitat de Catalunya under contract 2021-SGR-00478.

% \section*{Conflict of Interest}
% This work was partially funded by Lenovo and Intel\textsuperscript{\textregistered} through the Lenovo AI Innovators University Research program. The authors declare that this funding did not influence the judgements and actions of authors with regard to objective data presentation, analysis, and interpretation.

% \section*{Acknowledgements}
% We thank Lenovo for providing the technical infrastructure to run the experiments in this paper.

\section*{Competing Interests}
This work was partially funded by Lenovo and Intel\textsuperscript{\textregistered} through the Lenovo AI Innovators University Research program. The authors declare that this funding did not influence the judgements and actions of authors with regard to objective data presentation, analysis, and interpretation. The authors have no other relevant financial or non-financial interests to disclose.

\section*{Funding Information}
This work was partially supported by Lenovo and Intel\textsuperscript{\textregistered} as part of the Lenovo AI Innovators University Research program, by the Spanish Ministry of Science (MICINN), the Research State Agency (AEI) and European Regional Development Funds (ERDF/FEDER) under grant agreement PID2021-126248OB-I00, MCIN/AEI/10.13039/501100011033/FEDER, UE, and by the Generalitat de Catalunya under contract 2021-SGR-00478.

\section*{Author Contribution}
G. Bravo-Rocca conceived and designed the APE framework, conducted all experiments and data analysis, and wrote the manuscript. Co-authors provided scientific expertise, manuscript revision, and project oversight.

\subsection*{Data Availability Statement}
The APTOS 2019 Blindness Detection dataset used in this study is publicly available at Kaggle (\url{https://www.kaggle.com/c/aptos2019-blindness-detection}). The source code and experimental implementation for the APE framework are available at \url{https://github.com/gusseppe/ape}. 


\section*{Research Involving Human and/or Animals}
Not applicable.

\section*{Informed Consent}
Not applicable.

% \begin{appendices}

% \section{Section title of first appendix}\label{secA1}

% An appendix contains supplementary information that is not an essential part of the text itself but which may be helpful in providing a more comprehensive understanding of the research problem or it is information that is too cumbersome to be included in the body of the paper.

%%=============================================%%
%% For submissions to Nature Portfolio Journals %%
%% please use the heading ``Extended Data''.   %%
%%=============================================%%

%%=============================================================%%
%% Sample for another appendix section			       %%
%%=============================================================%%

%% \section{Example of another appendix section}\label{secA2}%
%% Appendices may be used for helpful, supporting or essential material that would otherwise 
%% clutter, break up or be distracting to the text. Appendices can consist of sections, figures, 
%% tables and equations etc.

% \end{appendices}

%%===========================================================================================%%
%% If you are submitting to one of the Nature Portfolio journals, using the eJP submission   %%
%% system, please include the references within the manuscript file itself. You may do this  %%
%% by copying the reference list from your .bbl file, paste it into the main manuscript .tex %%
%% file, and delete the associated \verb+\bibliography+ commands.                            %%
%%===========================================================================================%%

\bibliography{sn-bibliography}% common bib file
%% if required, the content of .bbl file can be included here once bbl is generated
%%\input sn-article.bbl

\end{document}
